\cotefile{158}
\scanpage{36}\mspage{16}

\[
\bar\varphi^{h}=1,
\]
donc $\bar\varphi$ est inversible. Mais alors $\deleted{\unsure{\bar g}}$
$\bar g\overset{\mathrm{déf}}{=}\bar\varphi^\beta$ est inversible (et $\unsure{=} $ d'ordre fini),
\deleted{Or dans un \unsure{monoïde}} et d'autre part il est minimal, comme
image de $\deleted{\unsure{u}}\varphi^\beta$ qui est minimal. \deleted{Ainsi,}
Mais si dans un monoïde $\bar M$ \deleted{tel que} \inserted{où}
on a la loi $[\bar u\bar f=\bar v\bar f\Longrightarrow\bar u=\bar v]$, on a
un élément minimal et inversible \inserted{$\unsure{\sigma|\varphi}$}, on
a vu que $\bar M$ est un groupe. Car
$\deleted{\unsure{\exists}}\forall u\in\bar M$, $\exists v\in M$ tel que $\deleted{\unsure{uv}}$
$\sigma=uv$ i.e. $u\underbrace{(v\sigma^{-1})}_{w}=1.$, et alors
$p=wu$ satisfait $p^2=p$ \deleted{d'où} $p.p=1.p$
i.e. $p=1$, \deleted{\unsure{donc}} $w$ est un inverse de
$u$, \deleted{\unsure{une}} tout élément est inversible.

On a prouvé :

\uline{Proposition 3} Soit $M$ un monoïde ayant
un \deleted{\unsure{ajo|objet}} élément minimal $\psi$, et soit
$\varphi\in M$ tel qu'il existe $\beta\in\mathbb N$ avec $\varphi^\beta$ minimal
(par exemple $\varphi$ minimal \deleted{\unsure{i.e. Si}}, \unsure{p. ex|p ce|par ex.}
$\varphi=\psi$). S'il existe $r,s\in\mathbb N$ avec $r\ne s$, et
$\varphi^r\psi=\varphi^s\psi$,

% Printed form upside down relative to the mathematics.
% letterhead
{\footnotesize
\begin{tabular}{@{}l@{\quad}l@{}}
\deleted{MINISTERE DE L'EDUCATION NATIONALE} & \deleted{AVIGNON, le}\\
\deleted{ACADEMIE D'AIX-MARSEILLE} &\\
\hfil\rule{1.5cm}{0.4pt}\hfil &\\
\deleted{Inspection Académique} & \deleted{L'Inspecteur d'académie,}\\
\deleted{de Vaucluse} & \deleted{directeur des services départementaux}\\
\hfil\rule{1.5cm}{0.4pt}\hfil & \deleted{de l'Education Nationale de Vaucluse,}\\
\deleted{Tél. : 90 82 99 76} &\\
\deleted{Référence : Secrétariat n\textsuperscript{o}} & \deleted{à M}
\end{tabular}}
% end letterhead

% archive sheet number: 35
